% 05_methodology
\section{\bgen{МЕТОДИКА НА ИЗСЛЕДВАНЕТО}{METHODOLOGY}}

Тази глава определя как се измерва, и я предхожда една констатация: измерването на
уеб сървъри е област, в която грешките рядко се проявяват като грешки. Неправилно
проведеното измерване не отказва да работи. То произвежда числа, които изглеждат
правдоподобни, подреждат се в графика и водят до заключение, изведено от нещо друго, а не
от това, което авторът смята, че е измерил.

Затова главата е подредена около начините, по които измерването може да бъде тихо
погрешно, и около механизмите, които правят всеки от тях невъзможен, а не невероятен.

\subsection{Какво трябва да произведе методиката}

Твърденията от глава~III определят изискванията, а не обратното. Те са три и всяко налага
различно ограничение.

Първо, цената на демултиплексирането се твърди, че е неоткриваема. Твърдение за липса на
разлика е по-трудно за защита от твърдение за наличие: то изисква разделителна способност,
достатъчна да покаже разликата, ако тя съществуваше.

Второ, делът на пакетите, изискващи препращане, се твърди, че е малък. Това изисква
условия, при които миграцията изобщо се случва, тоест повече от един адрес на клиента.

Трето, отложеното изобразяване се твърди, че намалява времето до първия байт. Това изисква
разделяне на времето до първия байт от времето до пълното изобразяване, тъй като второто
не се подобрява.

\subsection{Сравнени подходи}

Разглеждат се три подхода, преди да се приеме единият.

\paragraph{Измерване през loopback.} Най-простият: клиент и сървър на една машина, свързани
през \code{lo}. Отхвърля се, и не като ограничение, което да се спомене, а като
дисквалифициращо за централното твърдение.

Причината е, че loopback не е неутрален спрямо протокола. Интерфейсът \code{lo} има MTU
65536, тоест протоколите над TCP получават сегменти от 64\,KB без сегментиране, без
контролни суми и без прекъсвания, докато QUIC по замисъл продължава да изпраща дейтаграми
от 1200 до 1450 байта и плаща пълна криптография в потребителско пространство на всяка от
тях. Сравнение на HTTP/1.1 и HTTP/2 срещу HTTP/3 през loopback наказва HTTP/3 по причини,
които нямат нищо общо с която и да е от двете реализации.

Освен това времето за обръщане през loopback е от порядъка на десетки микросекунди, тоест
всяко свойство на мултиплексирането, свързано със скриване на закъснение, става невидимо. И
понеже адресът е един, миграцията на връзка е непроверима по определение.

Loopback остава полезен за микроизмервания на маршрутизатора, за цената на сериализацията и
за изрично обозначен „под на протоколните разходи“. За нищо друго.

\paragraph{Две физически машини.} Най-точният, и той е и целта. Изисква втора машина и
връзка, чиято пропускателна способност не се превръща сама в измерваната величина: при
1\,Gbit/s таванът е около 118\,MB/s, което прави всяко измерване с голям отговор измерване
на кабела.

\paragraph{Двойка мрежови пространства.} Приетият подход. Две мрежови пространства на една
машина, свързани с \code{veth}, с \code{netem} по всяка посока поотделно.

% Двете конфигурации за измерване, поставени една над друга.
%
% Рисунката трябва да носи едно твърдение преди да е прочетен ред текст: втората
% конфигурация допуска по-тесен кръг от изводи. Затова двата панела завършват с еднакво
% широка лента „обхват“. При основната конфигурация цялата лента е запълнена, при
% вторичната запълнени са само около две пети отляво, а останалото е празен пунктирен
% правоъгълник.
% Сравнението е геометрично и не зависи от четенето.
%
% Черно-бяла рисунка. Разграничението е по форма, по дебелина на линията и по сив тон:
% активните единици са gray!25, поясненията gray!10, ограниченията са в рамка с удебелен
% контур, а липсващата възможност е зачертана с кръст. Няма нито един цвят.
\begin{figure}[htbp]
\centering
\begin{tikzpicture}[
    % Вътрешните координати стигат до около 15,3 по хоризонтала. Мащабът 0,9 ги свива до
    % под 13,8 cm, тоест под ширината на текста 14,5 cm.
    scale=0.9, transform shape,
    % preamble.tex дефинира note два пъти: веднъж като чист текст и по-късно като
    % пунктирана сива кутия. Втората дефиниция е тази, която важи в дисертацията, тоест
    % без този ред всеки надпис тук се огражда с кутия, рамките се застъпват със
    % стрелките и с пунктирните пространства, а разликата между пояснение и ограничение
    % изчезва. Тук note се връща към чист текст, но само за тази рисунка.
    note/.style={align=left, font=\fontsize{8}{10}\selectfont},
    unit/.style={proc, minimum width=4.6cm, minimum height=0.95cm, fill=gray!25,
                 font=\fontsize{8}{9.5}\selectfont},
    small unit/.style={unit, minimum width=3.6cm},
    addr/.style={box, minimum width=4.6cm, minimum height=0.85cm, fill=gray!10,
                 font=\fontsize{7.5}{9}\selectfont},
    nslab/.style={note, align=center, font=\fontsize{8}{9.5}\selectfont\bfseries},
    prop/.style={box, text width=4.2cm, align=center, fill=gray!10, minimum height=1.05cm,
                 font=\fontsize{7.5}{9}\selectfont},
    lim/.style={box, very thick, text width=4.0cm, align=left, fill=gray!10,
                minimum height=1.70cm, font=\fontsize{7.5}{9}\selectfont},
    scopeok/.style={box, anchor=west, minimum height=1.25cm, fill=gray!25, align=center,
                    font=\fontsize{7.5}{9}\selectfont},
    scopeno/.style={box, anchor=west, minimum height=1.25cm, densely dashed, align=center,
                    font=\fontsize{7.5}{9}\selectfont},
    head/.style={align=center, font=\fontsize{9.5}{11.5}\selectfont\bfseries},
    sub/.style={note, align=center, text width=13.2cm, font=\fontsize{7.5}{9}\selectfont},
    tag/.style={note, align=center, font=\fontsize{7}{8.5}\selectfont, inner sep=1.5pt},
    panel/.style={draw, thick, rounded corners=3pt, inner sep=8pt},
  ]

  % =============================================================== основна конфигурация
  \node[head] (mt) at (7.70, 1.05) {Основна конфигурация (Linux): двойка мрежови
                                    пространства};
  \node[sub]  (ms) at (7.70, 0.45) {Две мрежови пространства на една машина, свързани с
                                    двойка \code{veth}.};

  % Пространството на генератора носи два адреса. Това не е украса: без втори адрес
  % миграцията на връзка няма как да се случи, тоест твърдението за нея е непроверимо.
  \node[nslab] (nsg) at (2.90, -0.30) {мрежово пространство \code{gen}};
  \node[unit]  (gen) at (2.90, -1.20) {генератор\\ h2load, wrk2, k6};
  \node[addr]  (adg) at (2.90, -2.50) {\code{veth-gen}: 10.0.0.2\\
                                       втори адрес: 10.0.0.3};
  \node[grp, fit=(nsg)(gen)(adg)] (gg) {};

  \node[nslab] (nss) at (12.50, -0.30) {мрежово пространство \code{srv}};
  \node[unit]  (srv) at (12.50, -1.20) {сървър\\ \code{coroute}};
  \node[addr]  (ads) at (12.50, -2.50) {\code{veth-srv}: 10.0.0.1};
  \node[grp, fit=(nss)(srv)(ads)] (gs) {};

  % Двете посоки се рисуват като две отделни стрелки, защото \code{netem} се задава на
  % изхода на всяко \code{veth} поотделно. Една двупосочна стрелка би подсказала точно
  % грешката, която разделът за loopback описва.
  \draw[flow] (gg.east |- gen) -- (gs.west |- gen);
  \node[tag] at (7.70, -0.68) {\code{netem} към сървъра\\
                               закъснение, загуба, \code{limit}};

  \node[box, fill=gray!10, font=\fontsize{7}{8.5}\selectfont] (vp) at (7.70, -1.88)
        {двойка \code{veth}};

  \draw[flow] (gs.west |- ads) -- (gg.east |- ads);
  \node[tag] at (7.70, -3.02) {\code{netem} към клиента\\ зададено поотделно};

  % Трите свойства, заради които тази конфигурация носи централното твърдение.
  \node[prop] (p1) at (2.90, -4.30) {истински MTU 1500 и истинска опашка \code{qdisc}};
  \node[prop] (p2) at (7.70, -4.30) {истински IP път, а не \code{lo}};
  \node[prop] (p3) at (12.50, -4.30) {два адреса на клиента, миграцията е проверима};

  \node[note, anchor=north west, align=left, text width=14.1cm,
        font=\fontsize{7.5}{9.5}\selectfont] (nt) at (0.62, -5.15)
       {Ядрата са изолирани, генераторът и сървърът не делят ядро.\\
        \code{netem} стои на изхода на всяко \code{veth} поотделно, тоест зададеното
        закъснение е закъснение за посока, а не за обръщане.\\
        Дължината на опашката \code{limit} се задава изрично, защото подразбиращите се
        1000 пакета превръщат зададената загуба в препълване.};

  \node[scopeok, text width=13.6cm, minimum width=14.15cm] (bar1) at (0.62, -7.60)
       {обхват: сравнение между HTTP/1.1, HTTP/2 и HTTP/3, миграция на QUIC, влошена
        мрежа, горни персентили};

  % Двете невидими точки държат двата панела с еднаква ширина. Без тях рамките се
  % получават с няколко милиметра разлика, което се чете като небрежност.
  \node[panel, fit={(mt)(ms)(gg)(gs)(p1)(p2)(p3)(nt)(bar1)(0.30,-1.20)(15.20,-1.20)}]
        (mpanel) {};

  % ============================================================= вторична конфигурация
  \node[head] (st) at (7.70, -9.20) {Вторична конфигурация (Windows): loopback};
  \node[sub]  (ss) at (7.70, -9.75) {Генераторът и сървърът са на един хост и говорят
                                     през \code{lo}.};

  % Кутиите тук са по-тесни от горните, за да остане място вдясно за това, което липсва.
  \node[nslab] (hl) at (6.10, -10.40) {един хост};
  \node[small unit] (gen2) at (2.40, -11.20) {генератор\\ h2load, k6};
  \node[box, fill=gray!10, font=\fontsize{8}{9.5}\selectfont] (lo) at (6.10, -11.20)
       {\code{lo}\\ 127.0.0.1};
  \node[small unit] (srv2) at (9.80, -11.20) {сървър\\ \code{coroute}};
  \draw[flow, {Stealth}-{Stealth}] (gen2) -- (lo);
  \draw[flow, {Stealth}-{Stealth}] (lo) -- (srv2);
  \node[grp, fit=(hl)(gen2)(lo)(srv2)] (gh) {};

  % Зачертаната кутия е единственият елемент с кръст в рисунката. Тя казва какво липсва,
  % и затова стои извън пунктираната рамка на хоста. Вътре в нея стои една дума, защото
  % кръстът минава през текста и по-дълъг надпис не би се четял.
  \node[box, densely dashed, minimum width=2.4cm, minimum height=0.9cm] (sh)
        at (13.40, -11.20) {};
  \draw[very thick] (sh.north west) -- (sh.south east);
  \draw[very thick] (sh.south west) -- (sh.north east);
  % Надписът се поставя след кръста и с бял фон, иначе линиите минават през буквите.
  \node[fill=white, inner sep=1.5pt, font=\fontsize{8}{9.5}\selectfont]
        at (sh) {\code{netem}};
  \node[tag, below=3pt of sh] (shl) {няма програмируемо\\
                                     оформяне на трафика,\\
                                     няма \code{tc}};

  % Ограниченията се назовават тук, а не се оставят за раздела с изводите.
  \node[lim] (l1) at (2.85, -13.75) {\textbf{MTU 65536.} TCP получава сегменти от
        64\,KB, а QUIC остава при 1200 до 1450 байта. Сравнението наказва HTTP/3.};
  \node[lim] (l2) at (7.70, -13.75) {\textbf{Обръщане в микросекунди.} Скриването на
        закъснение от мултиплексирането става невидимо.};
  \node[lim] (l3) at (12.55, -13.75) {\textbf{Един адрес.} Миграцията на връзка е
        непроверима по определение.};

  % Лентата има същата обща ширина като горната. Запълнената част е обхватът, празната е
  % това, което тази конфигурация не може да поддържа.
  \node[scopeok, text width=5.0cm, minimum width=5.60cm] (bar2) at (0.62, -15.55)
       {обхват: микроизмервания на маршрутизатора, цена на сериализацията, преносимост};
  \node[scopeno, text width=8.0cm, minimum width=8.55cm] (bar3) at (6.22, -15.55)
       {извън обхвата: сравнение между протоколите, миграция на връзка, влошена мрежа,
        горни персентили};

  \node[panel, fit={(st)(ss)(gh)(sh)(shl)(l1)(l2)(l3)(bar2)(bar3)
                    (0.30,-11.20)(15.20,-11.20)}] (spanel) {};

\end{tikzpicture}
\caption{Двете конфигурации за измерване. Горе: основната конфигурация под Linux, две
мрежови пространства на една машина, свързани с двойка \code{veth}, с \code{netem} на
изхода на всяка посока поотделно и с изрично зададена дължина на опашката. Тя дава
истински MTU, истинска опашка, истински IP път и два адреса на клиента, тоест миграцията
на QUIC е проверима. Долу: вторичната конфигурация под Windows, генератор и сървър на един
хост през \code{lo}, без програмируемо оформяне на трафика. Трите ограничения са
отбелязани изрично. Долните ленти имат еднаква обща ширина и сравняват обхвата на
твърденията, които всяка конфигурация допуска.}
\label{fig:measurement_topology}
\end{figure}


Той дава истински MTU, истинска опашка, истински IP път и два адреса на клиента, което е
единственото, което прави миграцията на QUIC изобщо проверима на една машина. Приема се
като основен, а измерванията на две машини се използват за проверка на валидността върху
малък брой клетки.

\paragraph{Уточнение за \code{netem}.} Прилагането на \code{netem} върху основната опашка
на \code{lo} е грешка, която си струва да се назове, защото не се проявява като грешка.
През loopback и заявката, и отговорът са изходящ трафик по един и същи интерфейс, тоест
зададено закъснение от 50\,ms доставя около 100\,ms за обръщане, а загуба от 1\% става
около 2\% на обръщане. Освен това \code{netem} по подразбиране има ограничение от 1000
пакета; при скоростите на loopback опашката прелива и „1\% загуба“ на практика е неизмерено
изчерпване на буфер.

\subsection{Инструменти за натоварване}

За сравнения между протоколи се използва \textbf{един} генератор: h2load, сглобен срещу
ngtcp2 и nghttp3.

Мотивът не е предпочитание, а елиминиране на променлива. h2load е единственият зрял
генератор, който говори HTTP/1.1, HTTP/2 и HTTP/3 с един и същ изходен формат и един и същ
модел на латентността. Сравнение между протоколи, направено с три различни клиента, е
объркано с три различни клиентски реализации, и няма как след това да се раздели кое на
какво се дължи. h2load освен това предоставя брой отговори по код на състоянието, което
затваря най-опасния пропуск, описан по-долу.

Две цени се обявяват вместо да се премълчат. Пътят за HTTP/1.1 на h2load е по-слабо
оптимизиран от този на wrk, тоест абсолютната пропускателна способност за HTTP/1.1 ще
изглежда по-ниска, отколкото при wrk. Това е приемливо, защото всяко сравнение се прави при
един клиент, но трябва да бъде казано, и една контролна клетка с wrk определя разликата
между клиентите, вместо тя да се предполага. Второ, ключът \code{-r} на h2load задава темп
на пристигане на \emph{връзки}, а не на заявки, тоест не замества wrk2.

Разпределението на задачите между три инструмента е следното: h2load за наситена
пропускателна способност и за разпределения на времето за обслужване по трите протокола;
wrk2 за HTTP/1.1 при постоянен темп на заявките в отворен цикъл; k6 за криви на латентността
спрямо предложеното натоварване и за натоварване на WebSocket.

\subsection{Съгласувано пропускане}

Затвореният цикъл измерва време за обслужване, а не време за отговор, и разликата не е
терминологична.

При затворен цикъл с $C$ връзки всяка връзка изпраща следващата заявка едва след като е
получила предходния отговор. Когато сървърът блокира, клиентът спира да изпраща. Заявките,
които \emph{биха} били забавени, никога не се изпращат и никога не се измерват. Полученият
99-и персентил е персентил на времето за обслужване на заявките, които са били допуснати, а
не на времето за отговор, каквото го вижда потребител. Двата цикъла са съпоставени на \figref{fig:closed_vs_open_loop}.

\begin{figure}[htbp]
\centering
\begin{tikzpicture}[
  x=1cm, y=1cm,
  % Двата панела делят една времева ос и една и съща сива лента за престоя.
  % Това е нарочно: сравнението има смисъл само ако престоят е един и същ.
  sbar/.style={draw, fill=gray!10, anchor=west, inner sep=0pt,
               minimum height=0.34cm},
  stalled/.style={sbar, fill=gray!40},
  % Един и същ знак за изпращане на заявка в двата панела. Броят му вътре в
  % лентата е цялото твърдение на фигурата, затова знакът трябва да е еднакъв.
  send/.style={line width=1pt},
  hdr/.style={anchor=west, font=\fontsize{9}{11}\selectfont},
  lane/.style={note, anchor=east, align=right},
  meas/.style={Stealth-Stealth, semithick},
  inbar/.style={font=\fontsize{7.5}{9}\selectfont},
]

% ----------------------------------------------------- престой на сървъра
% Лентата е прекъсната само там, където минава заглавието на долния панел.
% Двете парчета са на едни и същи x, за да се чете като един и същ престой.
\begin{scope}[on background layer]
  \fill[gray!12] (5.80,1.25) rectangle (8.00,-2.10);
  \fill[gray!12] (5.80,-2.78) rectangle (8.00,-6.00);
\end{scope}
\foreach \x in {5.80,8.00} {
  \draw[densely dashed, black!55] (\x,1.25) -- (\x,-2.10);
  \draw[densely dashed, black!55] (\x,-2.78) -- (\x,-6.00);
}
\node[note, anchor=south, align=center] at (6.90,1.32) {престой на сървъра};

% ===================================================== ЗАТВОРЕН ЦИКЪЛ
\node[hdr] at (0.35,1.95)
  {\textbf{Затворен цикъл} (\code{h2load}): връзката чака отговор преди следващата заявка.};

% Ред с изпращанията от трите връзки, слети на една ос.
\node[lane] at (1.85,1.02) {изпращания};
\foreach \x in {2.00,2.20,2.45,2.90,3.15,3.25,3.80,3.90,4.30,4.60,4.70,5.30,5.45,5.60}
  \draw[send] (\x,0.87) -- (\x,1.17);
\foreach \x in {8.15,8.20,8.25,8.90,9.15,9.25,9.65,10.05,10.30,10.40,10.95,11.15,11.35}
  \draw[send] (\x,0.87) -- (\x,1.17);

% Празнината между последното и първото следващо изпращане е цялата поанта.
\draw[meas] (5.60,0.70) -- (8.15,0.70);
\node[note, anchor=north, align=center] at (6.88,0.62) {нула изпращания};

% Три връзки. Всяка има точно една заявка в полет, когато престоят започне,
% затова всяка има точно една удължена лента и нищо друго вътре в престоя.
\node[lane] at (1.85,0.02) {връзка 1};
\foreach \x in {2.00,2.90,3.80,4.70} \node[sbar, minimum width=0.80cm] at (\x,0.02) {};
\node[stalled, minimum width=2.55cm] at (5.60,0.02) {};
\foreach \x in {8.25,9.15,10.05,10.95} \node[sbar, minimum width=0.80cm] at (\x,0.02) {};

\node[lane] at (1.85,-0.58) {връзка 2};
\foreach \x in {2.20,3.25,4.30} \node[sbar, minimum width=0.90cm] at (\x,-0.58) {};
\node[stalled, minimum width=2.65cm] at (5.45,-0.58) {};
\foreach \x in {8.20,9.25,10.30,11.35} \node[sbar, minimum width=0.90cm] at (\x,-0.58) {};

\node[lane] at (1.85,-1.18) {връзка 3};
\foreach \x in {2.45,3.15,3.90,4.60} \node[sbar, minimum width=0.60cm] at (\x,-1.18) {};
\node[stalled, minimum width=2.75cm] at (5.30,-1.18) {};
\foreach \x in {8.15,8.90,9.65,10.40,11.15} \node[sbar, minimum width=0.60cm] at (\x,-1.18) {};

% Мярката, която затвореният цикъл наистина дава.
\draw[densely dotted, gray!70] (5.30,-1.36) -- (5.30,-1.62);
\draw[densely dotted, gray!70] (8.05,-1.36) -- (8.05,-1.62);
\draw[meas] (5.30,-1.68) -- (8.05,-1.68);
\node[note, anchor=north, align=center, fill=white, inner sep=1.5pt] at (6.68,-1.76) {време за обслужване};

% ===================================================== ОТВОРЕН ЦИКЪЛ
\node[hdr] at (0.35,-2.50)
  {\textbf{Отворен цикъл} (\code{wrk2}): фиксиран темп на заявките, независимо от сървъра.};

\node[lane] at (1.85,-3.10) {заявки};
\draw[gray!55, thin] (1.95,-3.10) -- (12.30,-3.10);
\foreach \x in {2.00,2.43,...,12.30} \draw[send] (\x,-3.25) -- (\x,-2.95);

% Една проследена заявка. Пристига вътре в престоя, чака, после се обслужва.
\node[lane] at (1.85,-3.75) {една заявка};
\node[sbar, minimum width=2.60cm, inbar] at (6.30,-3.75) {чакане в опашката};
\node[stalled, minimum width=0.80cm] at (8.90,-3.75) {};
\node[note, anchor=west] at (9.80,-3.75) {обслужване};
\draw[densely dotted, gray!70] (6.30,-3.25) -- (6.30,-3.58);

\draw[meas] (6.30,-4.17) -- (9.70,-4.17);
\node[note, anchor=north, align=center, fill=white, inner sep=1.5pt] at (8.00,-4.25) {време за отговор};

% Дължина на опашката. Стълбата е стъпаловидна, защото пристиганията са
% дискретни: по едно качване на всяко пристигане, докато сървърът стои.
\node[lane] at (1.85,-5.30) {опашка};
\fill[gray!25]
  (2.00,-5.90) -- (5.87,-5.90) -- (5.87,-5.70) -- (6.30,-5.70) -- (6.30,-5.50) --
  (6.73,-5.50) -- (6.73,-5.30) -- (7.16,-5.30) -- (7.16,-5.10) -- (7.59,-5.10) --
  (7.59,-4.90) -- (8.02,-4.90) -- (8.02,-5.10) -- (8.45,-5.10) -- (8.45,-5.30) --
  (8.88,-5.30) -- (8.88,-5.50) -- (9.31,-5.50) -- (9.31,-5.70) -- (9.74,-5.70) --
  (9.74,-5.90) -- (12.30,-5.90) -- cycle;
\draw[thick]
  (2.00,-5.90) -- (5.87,-5.90) -- (5.87,-5.70) -- (6.30,-5.70) -- (6.30,-5.50) --
  (6.73,-5.50) -- (6.73,-5.30) -- (7.16,-5.30) -- (7.16,-5.10) -- (7.59,-5.10) --
  (7.59,-4.90) -- (8.02,-4.90) -- (8.02,-5.10) -- (8.45,-5.10) -- (8.45,-5.30) --
  (8.88,-5.30) -- (8.88,-5.50) -- (9.31,-5.50) -- (9.31,-5.70) -- (9.74,-5.70) --
  (9.74,-5.90) -- (12.30,-5.90);

% Основата на опашката служи и за времева ос на цялата фигура.
\draw[-Stealth, semithick] (1.95,-5.90) -- (12.55,-5.90);
\node[note, anchor=south east] at (12.55,-5.87) {време};

% ===================================================== следствие
\node[note, anchor=north west, text width=11.90cm] at (0.35,-6.30)
  {\textbf{Следствие.} p99 от затворен цикъл подценява опашката на разпределението.
   Клиентът е спрял да поражда точно онова натоварване, което би я произвело.};

\end{tikzpicture}
\caption{Съгласувано пропускане, показано като механизъм, а не като твърдение.
Горе всяка от $C$ връзки изпраща следваща заявка едва след получен отговор, затова
за целия престой на сървъра няма нито едно изпращане. Измерва се време за обслужване
на допуснатите заявки. Долу заявките се предлагат с постоянен темп, престоят не ги
спира, опашката расте и се източва след него. Измерва се време за отговор, тоест
чакане плюс обслужване. Затова 99-ият персентил от затворен цикъл подценява горния
край на разпределението: липсват именно заявките, които престоят би забавил най-много.}
\label{fig:closed_vs_open_loop}
\end{figure}


Протоколът е следният. Първо изкачване в затворен цикъл, за да се намери $T_{max}$ на всяка
система. След това $T_{ref}$ се приема за минимума от $T_{max}$ на сравняваните системи за
дадената клетка. Накрая се изпълняват измервания в отворен цикъл при фиксирани нива на
предложено натоварване от 25, 50, 75, 90 и 95 процента от $T_{ref}$, тоест всяка система
получава едно и също предложено натоварване.

Числата от затворен цикъл се съобщават изрично като време за обслужване, а тези от отворен
цикъл като време за отговор. Смесването на двете в една таблица е начинът, по който
разликата изчезва.

\paragraph{Заявен пропуск.} За HTTP/3 няма зрял генератор с постоянен темп на заявките в
отворен цикъл. wrk2 говори само HTTP/1.1, k6 и vegeta не говорят HTTP/3, а управлението на
темпа при h2load е по връзки. Следователно латентността в горните персентили за HTTP/3 се
съобщава само от затворен цикъл, изрично обозначена като време за обслужване, а
ограничението се заявява, вместо да се заобикаля.

\subsection{Едно рамо и друго рамо от един и същ двоичен файл}

Всяко сравнение между два режима на настоящата система се управлява с флаг при изпълнение,
а не с опция при компилиране. Това е ограничение върху реализацията, наложено от методиката,
и то е причината флаговете от глава~III да съществуват.

Мотивът е количествен. Разликите между отделни сборки, дължащи се на подредбата на кода, на
реда на свързване и на размера на средата, са документирани в порядъка от пет до десет
процента \cite{mytkowicz2009wrong}, докато разликите между изпълнения на една и съща
контролирана машина са порядъчно по-малки. Следователно разлика от три процента между два
поотделно компилирани двоични файла е неразличима от това обектните файлове да са били
свързани в друг ред.

Едно и също изпълнимо, два флага. Проверява се, че двете рамена наистина произхождат от
един файл, чрез сравнение на контролната му сума.

\subsection{Средата като контролирана променлива}

Кампанията продължава седмици. През това време ядрото се обновява, управлението на честотата
се връща на пестелив режим след приспиване, зависимост се преизгражда. Всяко от тези
събития прави изпълненията преди и след него несравними, и нито едно от тях не се обявява.
Числата остават правдоподобни, което е и проблемът.

Затова средата се записва и от нея се извлича отпечатък, а програмата, която управлява
кампанията, отказва да добави изпълнения към кампания, чийто отпечатък се е променил.
Продължаването въпреки това изисква изричен ключ, тоест решение, което някой е взел, а не
случайност.

Интересната преценка не е как се хешира, а какво се хешира. Записва се всичко; в отпечатъка
влиза само това, чиято промяна прави сравнението невалидно: ядро, модел на процесора, брой
физически и логически ядра, режим на управление на честотата, настройки за големи страници,
компилатор, версия на кода и версии на зависимостите. Средното натоварване и времето от
стартирането се записват и \emph{не} влизат, защото хеширане на всичко би направило всяко
изпълнение отделна кампания, а предпазна мярка, която се задейства постоянно, е предпазна
мярка, която някой изключва.

Отказът назовава полето, което се е променило. Съобщение „средата се промени“ се
пренебрегва от раздразнение; съобщение „режимът на честотата премина от \code{performance}
към \code{powersave}“ води до действие.

\subsection{Правила за допустимост, обявени предварително}

Всеки критерий за отхвърляне е определен преди събирането на данни, и този ред е същината.
Критерии, избрани след като резултатите са видени, са неразличими от избиране на резултати:
винаги съществува защитимо на пръв поглед основание да отпадне точно изпълнението, което
разваля тенденцията.

Критериите са малко на брой и всеки от тях назовава конкретна механична повреда, а не
мнение дали числото изглежда правилно.

\begin{table}[H]
\centering
\caption{Предварително обявени критерии за отхвърляне на изпълнение}
\label{tab:validity}
\begin{tabular}{@{}L{5.2cm}L{3.0cm}L{5.4cm}@{}}
\toprule
\textbf{Условие} & \textbf{Праг} & \textbf{Какво означава} \\
\midrule
Дял отговори извън 2xx & над 0,1\% & Сървърът е отказвал, а не обслужвал. \\
Натоварване на генератора & над 85\% & Измерва се клиентът, не сървърът. \\
Отклонение на честотата & над 2\% & Дроселиране по температура. \\
\code{UdpRcvbufErrors}, \code{TcpExtListenOverflows} & всяко нарастване & Ядрото е изпуснало работа, преди сървърът да я види. \\
Виртуализация & открита & Планирането и таймерите не са на машината. \\
Работно дърво с некоммитнати промени & открито & Записаната версия не описва двоичния файл. \\
\bottomrule
\end{tabular}
\end{table}

Първият критерий е най-важният. Без него сървър, който се срива под натоварване, печели по
пропускателна способност, защото отказът е по-евтин от обслужването.

Обратното също е част от правилото: \textbf{бавното изпълнение се приема}. Бавно е резултат,
и критерий, който отхвърля бавни изпълнения, отхвърля находката.

Отхвърлените изпълнения се запазват, а броят им се съобщава заедно с резултатите. Методика,
която е отхвърлила една трета от изпълненията си, е различна методика от такава, която не е
отхвърлила нито едно, и само съобщаването прави разликата видима.

\subsection{Ред на изпълнение}

Очевидният ред е всички изпълнения на система A, след това всички на система B. Той е и
единственият ред, гарантирано погрешен: машината се загрява, и тогава всяка разлика между
първия и втория час се приписва на разликата между A и B.

Затова планът се обхожда на проходи. Всеки проход посещава всяка клетка по веднъж, системите
се редуват, а редът вътре в прохода е разбъркан, при това различно за всеки проход.
Загряването продължава да съществува, но се разпределя приблизително по равно, а индексът на
прохода го прави видимо: ако пети проход е по-бавен от първи за всичко, това е температура и
може да се види, вместо да се извежда.

Разбъркването е с фиксирано семе. Кампания, която не може да бъде повторена в същия ред, не
е възпроизводима, а „случаен ред“ не е основание за неповторимост.

\subsection{Статистика}

Съобщават се медиани, а не средни стойности. Резултатите на ниво изпълнение не са
симетрични: едно температурно неудачно изпълнение измества средната стойност и почти не
помръдва медианата. Латентните извадки \textbf{не се подрязват} при никакви обстоятелства,
тъй като горните персентили са самото явление; следователно обобщаващата статистика трябва да
е устойчивата.

Доверителните интервали се получават чрез бутстрап с корекция за отклонение и ускорение
\cite{efron1987better}. Процентилният бутстрап е по-кратък за реализиране и е грешен по
начин, който тук има значение: той предполага, че бутстрап разпределението е без отклонение
и с постоянна дисперсия, а за медиана от пет или десет изпълнения то не е нито едното.

Разлика между системи се съобщава като такава само когато изпълни \textbf{две} условия:
доверителният интервал на разликата изключва нулата, \emph{и} относителната разлика е поне
пет процента. Статистическата различимост не е същото като значимост: при достатъчно
изпълнения разлика от 0,4\% става значима и остава безинтересна.

Броят на повторенията не се обявява, а се определя: една клетка се изпълнява двадесет пъти,
за да се измери променливостта между изпълнения, и от нея се избира $n$.

\subsection{\bgen{Изводи}{Summary}}

Методиката се свежда до няколко механизма, всеки от които прави определен вид тиха грешка
невъзможна, вместо да я прави малко вероятна.

Двойката мрежови пространства прави сравнението между протоколите честно, което loopback не
може. Един генератор прави сравнението между протоколите сравнение между протоколи, а не
между клиенти. Отвореният цикъл разделя времето за отговор от времето за обслужване. Един
двоичен файл с флагове прави разликата между две рамена по-голяма от шума между сборки.
Отпечатъкът на средата спира кампанията, вместо да смеси две популации. Предварително
обявените критерии премахват свободата резултатите да бъдат избирани след като са видени.

Три ограничения се заявяват тук, а не се откриват по-късно. За HTTP/3 няма зрял генератор с
постоянен темп на заявките в отворен цикъл, тоест горните персентили за него идват само от
затворен цикъл. Под Windows няма програмируем еквивалент на \code{tc}, тоест измерванията
при влошена мрежа не се извършват там, а разликите в UDP сегментирането правят твърденията
за производителност на HTTP/3 валидни само за Linux. И измерванията на две машини се
ограничават до малък брой клетки, използвани за проверка на валидността, а не за отделна
матрица.
